\chapter{Refrigerator and Heat Pump}\label{ch:refrig}

The Joule expansion of the previous chapter is the working stroke of both a
\emph{refrigerator} (harvesting the cooled left chamber) and a \emph{heat pump}
(harvesting the warmed right chamber): one and the same device, distinguished
only by which reservoir's heat exchange is counted as useful. Because RNS
resolves the dynamics, we can read their efficiency directly off the computed
fields rather than postulate it.

\section{A Computable Model}

An interactive realisation (\texttt{joule\_experiment\_cpu.html}, a
dependency-free vanilla-JS port of the original sketch) solves the 2D
compressible Navier--Stokes/RNS system
\[
\dot\rho+\nabla\cdot m=0,\qquad
\dot m+\nabla\cdot(mu)+\gamma\nabla e=\mu_{\mathrm{mom}}\,\Delta u,\qquad
\dot e+\nabla\cdot(eu)+\gamma e\,\nabla\cdot u=\mu_{\mathrm{mom}}\,|\nabla u|^2,
\]
with $p=\gamma e=\gamma\rho T$, $T=e/\rho$, $u=m/\rho$, on two chambers joined by a
channel: a dense, equally-warm left chamber ($\rho=e=10$, hence $T=1$) expanding
into a dilute right chamber ($\rho=e=1$). The last term in the energy equation is
the \emph{turbulent dissipation} $\mu_{\mathrm{mom}}|\nabla u|^2\ge 0$ --- the
irreversible return of the jet's kinetic energy to heat --- and crucially it uses
the \emph{same} coefficient $\mu_{\mathrm{mom}}$ as the momentum diffusion, so the
kinetic energy removed by viscosity reappears exactly as heat in the energy
equation and $K+E$ is conserved. A separate, much smaller coefficient
$\mu\ll\mu_{\mathrm{mom}}$ provides numerical stabilisation of $\rho$ and $e$
through Laplacian smoothing; it is not a physical dissipation and contributes
negligibly to the energy budget.

\section{Energy and Dissipation Bookkeeping}

The simulation tracks, at each instant, the kinetic energy
$K=\int\tfrac12\rho|u|^2\,dx$ and the internal energy $E=\int e\,dx$, with the
total $K+E$ conserved (Joule's null result for the \emph{whole} box), together
with the \emph{cumulative turbulent dissipation}
\begin{equation}\label{D_def}
D(t)=\int_0^t\!\!\int \mu_{\mathrm{mom}}\,|\nabla u|^2\,dx\,d\tau \;\;>\;0 ,
\end{equation}
which quantifies the irreversibility directly: it is positive by construction of
the discrete scheme and is \emph{the} 2nd Law for this model. There is no need
to introduce a statistical entropy $S$ or a relation $dS\ge 0$ --- the inequality
$D>0$ is the 2nd Law in its native, deterministic form, deriving from finite-
precision computation rather than postulated from probabilistic foundations.
$D$ accumulates wherever the velocity gradient is large: in the channel jet and
in the recompression against the far wall.

\section{Refrigerator and Heat Pump: Efficiency and the Cost of the Cycle}

A refrigerator runs \emph{cyclically}: each cycle the working gas must be
restored to its initial compressed state, so the relevant energy cost is the
work to restore the initial conditions. Its reversible (minimal) value is the
\emph{exergy} of the initial state; for the isothermal initial imbalance (both
chambers start at $T=T_0=1$) this is
\begin{equation}\label{exergyA0}
A_0=\gamma\int \rho\,\ln(\rho/\bar\rho)\,dx \;=\; D_\infty ,
\end{equation}
with $\bar\rho$ the mean density; the second equality states that the spontaneous
free expansion eventually destroys exactly this $A_0$ as accumulated turbulent
dissipation $D_\infty$ (\ref{D_def}). The cooling delivered is the heat removed from the cells
that drop below ambient,
\[
Q_{\mathrm{cold}}=\int_{T<T_0}\rho\,(T_0-T)\,dx ,
\]
and the coefficient of performance --- the buyer's figure of merit, cooling per
unit of work spent restoring the cycle --- is
\begin{equation}\label{copjoule}
\mathrm{COP}_{\mathrm R}=\frac{Q_{\mathrm{cold}}}{A_0}.
\end{equation}
The \emph{same} expansion is a \emph{heat pump} if instead we harvest the warmed
right chamber. Its useful output is the heat delivered to the hot side, which
over the full cycle is $Q_{\mathrm{hot}}=Q_{\mathrm{cold}}+A_0$ --- the cooling
moved up, plus the restore work $A_0$ rejected as heat during recompression ---
so
\[
\mathrm{COP}_{\mathrm{HP}}=\frac{Q_{\mathrm{cold}}+A_0}{A_0}=\mathrm{COP}_{\mathrm R}+1 .
\]
Refrigerator and heat pump are thus one device, the single Joule expansion,
distinguished only by which reservoir's heat exchange is counted as useful; in
the expansion phase alone the simulation shows $Q_{\mathrm{hot}}\approx
Q_{\mathrm{cold}}$ (pure redistribution at conserved total energy), the extra
$A_0$ entering on the hot side only during recompression.

Against the Carnot ceiling for the achieved lift,
$\mathrm{COP}_{\mathrm{Carnot}}=T_{\mathrm{cold}}/(T_{\mathrm{hot}}-T_{\mathrm{cold}})$,
the second-law efficiency $\eta=\mathrm{COP}_{\mathrm R}/\mathrm{COP}_{\mathrm{Carnot}}$ comes
out at only a few percent: the free expansion is a \emph{very poor} refrigerator,
destroying almost all of the available work $A_0$ as cumulative turbulent
dissipation $D$ while delivering only a sliver of cold. This is exactly why a
practical refrigerator does not free-expand but runs the recompression cycle of
the next chapter, recovering work and driving the lost-work $D$ down toward zero
(the Carnot limit). The same bookkeeping --- $A_0$, $Q_{\mathrm{cold}}$, $D$ ---
will let us assign an efficiency to any continuum-thermodynamic process,
including the thermodynamics of the living cell.

